% !TeX program = xelatex
\documentclass[12pt,a4paper]{article}
\usepackage{amsmath,amssymb}
\usepackage{geometry}
\usepackage{booktabs}
\usepackage{graphicx}
\usepackage{xcolor}
\usepackage[UTF8]{ctex}
\geometry{margin=2.5cm}
\title{Q1: 基础 Logistic 模型（零自由参数）}
\author{}
\date{}

\begin{document}
\maketitle

\section{模型建立}

基于经典 Logistic 增长理论，考虑到回收系统在实际运行过程中存在资源
容量约束与增长速率限制，假设回收量 $x(t)$ 随时间呈现先快速增长、
后逐渐趋缓并最终趋于稳定的动态演化规律。其中，固有增长率 $\gamma$
 描述系统在初始阶段的增长潜力，环境容量 $K$ 表示在现有条件下系统能
 够达到的最大稳定回收水平。由于本节主要用于建立基础预测模型，暂不考
 虑外部促进因素和内部抑制机制，仅保留两个核心约束因素：

\begin{equation}\label{eq:ode}
\frac{dx}{dt} = \gamma \cdot x \cdot \left(1 - \frac{x}{K}\right), \qquad x(0) = x_0
\end{equation}

解析解：

\begin{equation}\label{eq:solution}
x(t) = \frac{K}{1 + \left(\dfrac{K}{x_0} - 1\right) e^{-\gamma t}}
     = \frac{K}{1 + A \cdot e^{-\gamma t}}, \qquad A = \frac{K}{x_0} - 1
\end{equation}

\textbf{注意}：由于本模型中的参数 $\gamma$、$K$ 和初始状态 
$x_0$ 均来源于已知条件，因此该模型属于无参数调整的确定性模型。
模型预测结果完全由先验参数决定，不通过数据拟合进一步优化参数，
从而能够用于检验给定理论增长机制与实际恢复过程之间的一致性。
\section{参数}

\begin{center}
\begin{tabular}{lrll}
\toprule
参数 & 符号 & 数值 & 单位 \\
\midrule
饱和容量      & $K$      & 8011   & tons/day \\
固有增长率    & $\gamma$ & 0.022  & 1/day \\
初始回收量    & $x_0$    & 6314   & tons/day \\
导出常数      & $A$      & 0.2688 & --- \\
\bottomrule
\end{tabular}
\end{center}

\section{模型评估}

为定量分析基础 Logistic 模型对实际恢复过程的描述能力，
本文选取决定系数 $R^2$、平均绝对误差（MAE）以及平均绝对百分比误差（MAPE）
作为模型评价指标。

其中，$R^2$ 用于衡量模型对数据整体变化趋势的解释能力，
MAE 和 MAPE 分别从绝对误差和相对误差角度评价模型预测结果的偏差程度。

\begin{center}
\begin{tabular}{lc}
\toprule
指标 & 数值 \\
\midrule
$R^2$  & 0.5319 \\
MAE    & 315.6 \\
MAPE   & 4.42\% \\
\bottomrule
\end{tabular}
\end{center}

\section{预测与残差}

\begin{center}
\begin{tabular}{crrrr}
\toprule
$t$ (days) & Actual & Predicted & Residual & Rel.\ Error \\
\midrule
0   & 6314 & 6314.0 & $-0.0$   & $-0.00\%$ \\
30  & 6542 & 7033.9 & $-491.9$ & $-7.52\%$ \\
60  & 6875 & 7474.4 & $-599.4$ & $-8.72\%$ \\
90  & 7173 & 7724.4 & $-551.4$ & $-7.69\%$ \\
120 & 7368 & 7860.2 & $-492.2$ & $-6.68\%$ \\
150 & 7591 & 7932.4 & $-341.4$ & $-4.50\%$ \\
180 & 7724 & 7970.2 & $-246.2$ & $-3.19\%$ \\
270 & 7896 & 8005.3 & $-109.3$ & $-1.38\%$ \\
365 & 8002 & 8010.3 & $-8.3$   & $-0.10\%$ \\
\bottomrule
\end{tabular}
\end{center}

\textbf{系统性偏差}：从残差分布可以发现，预测值在整个时间区间内均高于实际观测值，
说明基础 Logistic 模型存在一定程度的系统性高估现象。

其中，最大负残差出现在 $t=60$ 天，对应偏差为
$-599.4$ tons/day。

该结果表明，在恢复初期阶段，实际系统增长速度低于理论模型预设水平。
产生该偏差的原因可能包括资源配置限制、运行效率波动以及外部环境因素影响，
因此固定增长率 $\gamma=0.022$ 难以完全反映真实恢复过程中的动态变化。

\section{远期预测}

\begin{center}
\begin{tabular}{ccc}
\toprule
$t$ (days) & $x(t)$ (tons/day) & 占 $K$ 比例 \\
\midrule
400 & 8010.7 & 99.996\% \\
450 & 8010.9 & 99.999\% \\
500 & 8011.0 & 100.00\% \\
600 & 8011.0 & 100.00\% \\
730 & 8011.0 & 100.00\% \\
\bottomrule
\end{tabular}
\end{center}

根据预测结果，系统在约400天后逐渐进入稳定阶段，
回收量持续趋近理论环境容量 $K=8011$ tons/day。

由于基础 Logistic 模型假设环境容量保持恒定，
因此长期预测表现为稳定的渐近收敛过程，
不会出现超过容量限制的增长现象。

该结果体现了容量约束机制对于长期恢复过程的主导作用。

\section{关键结论}

\begin{enumerate}

\item
\textbf{基础模型解释能力有限：}
模型获得决定系数 $R^2=0.532$，表明固定参数 Logistic 模型能够描述恢复过程的总体增长趋势，
但对于中期阶段的动态变化刻画能力不足。

特别是在 $30$--$150$ 天时间范围内，模型预测值持续高于实际观测值，
说明基础模型未充分考虑恢复过程中的阶段性影响因素。

\item
\textbf{增长机制存在动态偏差：}
残差分析显示模型预测增长速度整体偏快，
说明固定增长率 $\gamma=0.022$ 难以完全反映实际恢复过程中的复杂变化。

因此，在后续模型优化中，可通过引入抑制参数 $\beta$
对有效增长率进行动态修正，提高模型对实际演化规律的适应能力。

\item
\textbf{环境容量假设存在修正空间：}
基础 Logistic 模型假设环境容量 $K$ 为固定常数，
而实际系统可能受到技术提升、资源投入以及政策支持等外部因素影响，
使长期稳定水平发生变化。

因此，后续模型可以考虑引入动态容量参数 $K_{eff}$，
增强模型对外部促进因素的描述能力。

\item
\textbf{后续模型改进方向：}
针对基础模型存在的系统性偏差，
下一阶段将在 Q1 模型基础上引入促进系数 $\alpha$
和抑制系数 $\beta$，
分别调整有效容量 $K_{eff}$ 和有效增长率 $r$，
从而建立更加符合实际恢复过程的扩展模型。

\end{enumerate}

\end{document}
